% !TEX program = xelatex
% !TEX encoding = UTF-8 Unicode
\RequirePackage{iftex}
\ifPDFTeX
  \documentclass[UTF8,11pt,a4paper,fontset=none]{ctexart}
  \renewcommand{\CJKrmdefault}{gbsn}
  \renewcommand{\CJKsfdefault}{gbsn}
  \renewcommand{\CJKttdefault}{gbsn}
\else
  \documentclass[UTF8,11pt,a4paper,fontset=fandol]{ctexart}
  \usepackage{fontspec}
\fi
\usepackage[margin=22mm,headheight=14pt]{geometry}
\usepackage{amsmath,amssymb,xcolor,listings,fvextra}
\usepackage{array,booktabs,needspace,fancyhdr,graphicx}
\usepackage[unicode,hidelinks]{hyperref}
\definecolor{codebg}{RGB}{247,248,250}
\definecolor{coderule}{RGB}{220,224,229}
\definecolor{codekey}{RGB}{30,65,100}
\setlength{\parindent}{2em}
\setlength{\parskip}{6pt}
\setlength{\emergencystretch}{2em}
\linespread{1.15}
\ctexset{section={afterindent=true,format=\Large\bfseries,beforeskip=18pt,afterskip=9pt},subsection={afterindent=true,format=\large\bfseries,beforeskip=14pt,afterskip=7pt}}
\setlength{\abovedisplayskip}{9pt plus 2pt minus 2pt}
\setlength{\belowdisplayskip}{9pt plus 2pt minus 2pt}
\setlength{\jot}{4pt}
\newcommand{\uval}[1]{\left[#1\right]_{32}}
\newcommand{\hex}[1]{\mathtt{0x#1}}
\newcommand{\band}{\mathbin{\&}}
\newcommand{\shr}[2]{\left\lfloor\dfrac{#1}{2^{#2}}\right\rfloor}
\pagestyle{fancy}
\fancyhf{}
\fancyhead[L]{\small BombLab 实验报告}
\fancyhead[R]{\small 俞凡}
\fancyfoot[C]{\small\thepage}
\renewcommand{\headrulewidth}{0.3pt}
\lstset{language=C,basicstyle=\ttfamily\fontsize{9}{11}\selectfont,keywordstyle=\color{codekey},commentstyle=\color{gray},stringstyle=\color{codekey},backgroundcolor=\color{codebg},frame=single,rulecolor=\color{coderule},framerule=0.3pt,framesep=5pt,xleftmargin=6pt,xrightmargin=6pt,showstringspaces=false,columns=fullflexible,keepspaces=true,tabsize=4,breaklines=true,breakatwhitespace=false,aboveskip=6pt,belowskip=8pt}
\newcommand{\reportimage}[3]{%
  \par\addvspace{5pt}\noindent\begin{minipage}{\linewidth}\centering
  \includegraphics[width=0.96\linewidth,height=#2\textheight,keepaspectratio]{#1}\par
  \vspace{3pt}{\small #3}\par
  \end{minipage}\par\addvspace{6pt}}
\hypersetup{pdftitle={BombLab 实验报告},pdfauthor={俞凡}}
\begin{document}
\begin{center}{\LARGE\bfseries BombLab 实验报告}\end{center}


\noindent\textbf{姓名：俞凡\quad 学号：25301980234}\par

完整实现见 \href{run:defuser.c}{defuser.c}。我实际运行了单关测试、额外 seed 验证和总测试，截图放在对应关卡及结果汇总中。

\Needspace{6\baselineskip}\section*{一、实验内容与环境}


本实验在 VMware 的 Ubuntu x86-64 虚拟机中进行，通过 Windows 的 VS Code Remote SSH 编辑和调试。实验目录为 \texttt{/home/yufan/ICS/BombLab}，使用 GCC、Make、GDB 及课程提供的公开测试脚本。当前代码只在 \texttt{defuser.c} 中实现求解逻辑，个人炸弹、驱动、接口和测试程序保持原样。

实验的输入不是一组固定口令。驱动根据 seed 和关卡生成 64 位的 \texttt{challenge}，求解函数接收它并计算口令，再由使用相同参数的 \texttt{bomb} 验证。不同 seed 会生成不同题目，因此需要恢复通用计算关系，不能只保存某次运行得到的答案。

分析时先看输入检查，确定整数个数，再沿调用和比较寻找每个输出的来源。下面只选关键汇编，不逐行翻译整份程序。AT\&T 语法中源操作数在前、目的操作数在后；参数和返回值要结合当前调用观察，不能假定一个寄存器在整个函数里始终保存同一个变量。

以下公式中，$c$ 表示 challenge，$h$ 表示 context，并记
\[
  \uval{x}=x\bmod 2^{32},\qquad 0\leq \uval{x}<2^{32}.
\]
符号 $\oplus$ 表示逐位异或，$\band$ 表示逐位与；$x\gg n$ 表示无符号右移，即 $\lfloor x/2^n\rfloor$。除特别标明外，中间量都对应代码中的 \texttt{uint32\_t}，加法、乘法和左移的结果均保留低 32 位。这样既与汇编写入 \texttt{eax} 等 32 位寄存器的行为一致，也避免依赖 C 的有符号溢出。

\Needspace{6\baselineskip}\section*{二、逐关分析与实现}


\Needspace{189pt}\subsection*{第 1 关：跟踪 helper 的参数和返回值}


\texttt{phase\_01} 的输入检查要求一个整数。它将 challenge 交给 \texttt{p1\_compute}，再把返回的 \texttt{eax} 与这个整数比较。本关的重点是跟清 helper 的实际参数，而不是根据函数名字猜调用顺序。

以下指令分别来自三个 helper 和调用点：

\noindent\begin{minipage}{\linewidth}
\begin{lstlisting}[language={}]
402024: lea (%rdi,%rsi,4),%eax
402027: xor $0x5a0,%eax
402034: lea 0x46fa(%rsi),%eax
40203a: xor %edi,%eax
402044: xor $0xfb73,%esi
40204a: lea (%rsi,%rdi,8),%eax
40206d: call 402040 <p1_step2>
40207a: call 402020 <p1_step0>
402083: call 402030 <p1_step1>
\end{lstlisting}
\end{minipage}


三个 helper 对应的计算关系分别为
\[
\begin{aligned}
f_0(x,y)&=\uval{x+4y}\oplus\hex{5a0},\\
f_1(x,y)&=\uval{y+\hex{46fa}}\oplus x,\\
f_2(x,y)&=\uval{(y\oplus\hex{fb73})+8x}.
\end{aligned}
\]
\texttt{p1\_compute} 先以 high、low 为参数调用 step2，再以 low、mid 调用 step0，最后把 a、b 交给 step1。令 $c_\ell$、$c_m$、$c_h$ 分别对应代码中的 low、mid、high，则最终口令为
\[
\begin{aligned}
c_\ell&=\uval{c},\qquad c_m=\uval{c\gg16},\qquad c_h=\uval{c\gg32},\\[3pt]
a&=\uval{c_\ell+4c_m}\oplus\hex{5a0},\\
b&=\uval{(c_\ell\oplus\hex{fb73})+8c_h},\\
v&=a\oplus\uval{b+\hex{46fa}}.
\end{aligned}
\]


这里的 \texttt{mid} 与 low、high 有重叠，并不是把 64 位数分成三个互不相交的部分。\texttt{lea} 在这几处只进行整数表达式计算，不读取所算地址的内容。实现直接保留 low、mid、high、a、b，最后用 \texttt{write\_answer} 输出一个数，便于逐项与汇编核对。

使用 seed \texttt{42} 运行第 1 关，终端输出 \texttt{Phase 1 defused.}，随后执行单关公开测试，10 个 seed 全部通过。这里的 seed \texttt{42} 属于公开列表。

\reportimage{images/phase1-seed42.png}{0.47}{图 1：运行第 1 关，seed 为 42，程序输出 \texttt{Phase 1 defused.}}


\reportimage{images/phase1-public.png}{0.34}{图 2：第 1 关公开测试通过 10/10。}

\reportimage{images/phase1-seed8.png}{0.45}{图 3：第 1 关使用非公开 seed 8 运行，输出 \texttt{Phase 1 defused.}}


\Needspace{200pt}\subsection*{第 2 关：从栈上数组恢复递推关系}


输入检查和 \texttt{sscanf} 的返回值比较都要求六个整数。六个槽从 \texttt{rbp-0x40} 开始，相邻地址差 4 字节，与 32 位整数数组一致。第一项先从 challenge 取位，后面五项在循环中验证。

\noindent\begin{minipage}{\linewidth}
\begin{lstlisting}[language={}]
402160: shr $0x5,%rax
402164: and $0x1f,%eax
402167: add $0x10,%eax
40217b: shr %cl,%rax
40217e: and $0x1f,%eax
402181: xor %edx,%eax
402183: lea (%rax,%rsi,2),%eax
402186: mov -0x40(%rbp,%rdx,4),%esi
40218a: cmp %eax,%esi
402192: add $0x2,%ecx
\end{lstlisting}
\end{minipage}


循环开始时 $i=1$、移位量为 $10$；每轮 $i$ 增加 $1$，移位量增加 $2$，所以移位量为 $8+2i$。\texttt{lea} 执行时，\texttt{esi} 还是上一项，随后程序才从数组读取当前项并比较。六个整数满足
\[
\begin{aligned}
v_0&=((c\gg5)\band31)+16,\\
v_i&=\uval{2v_{i-1}+\bigl(((c\gg(8+2i))\band31)\oplus i\bigr)},
\quad i=1,\ldots,5.
\end{aligned}
\]


实现用六元素数组和一个循环，逐项生成，不把六项写成重复代码。判断顺序很重要：如果把读取当前项误当成发生在 \texttt{lea} 之前，就会把递推关系读反。循环结束条件也表明总共生成六项，而不是再额外计算第七项。

第 2 关公开测试通过 10/10。又使用公开列表之外的 seed \texttt{8} 单独运行第 2 关，输出 \texttt{Phase 2 defused.}，截图中的 seed 是 \texttt{8}。

\reportimage{images/phase2-public.png}{0.34}{图 4：第 2 关公开测试通过 10/10。}


\reportimage{images/phase2-seed8.png}{0.47}{图 5：使用非公开 seed 8 验证第 2 关，程序输出 \texttt{Phase 2 defused.}}


\Needspace{189pt}\subsection*{第 3 关：按跳转表恢复八个分支}


本关需要两个整数。第一个整数由 challenge 的低三位决定；第二个整数则由同一个索引选择计算分支。入口先构造传给分派函数的 x，分派函数再从地址表跳到实际代码。

\noindent\begin{minipage}{\linewidth}
\begin{lstlisting}[language={}]
4022f1: and $0x7,%edi
4022f4: cmp %edi,-0x18(%rbp)
4022ff: shr $0x3,%rax
402303: shr $0x20,%rsi
402307: xor %eax,%esi
402309: call 402240 <p3_dispatch>
402248: cmp $0x7,%eax
40224b: ja 402292 <p3_dispatch+0x52>
40224d: notrack jmp *0x405820(,%rax,8)
\end{lstlisting}
\end{minipage}


由入口指令可得分支索引和参与计算的数：
\[
  i=c\band7,\qquad x=\uval{c\gg3}\oplus\uval{c\gg32}.
\]
表项占 8 字节，存的是代码地址；应先读表内目标，再追踪该地址处的跳转。实际索引顺序与 \texttt{p3\_case\_N} 的编号不同，直接照函数名顺序填写会出错。

\par\addvspace{5pt}\noindent\begin{minipage}{\linewidth}\centering\small\renewcommand{\arraystretch}{1.28}
\begin{tabular}{>{\raggedright\arraybackslash}p{0.130\dimexpr\linewidth-6\tabcolsep\relax}>{\raggedright\arraybackslash}p{0.200\dimexpr\linewidth-6\tabcolsep\relax}>{\raggedright\arraybackslash}p{0.670\dimexpr\linewidth-6\tabcolsep\relax}}
\toprule
\textbf{索引 $i$} & \textbf{实际分支} & \textbf{第二个整数} \\
\midrule
0 & case 5 & $\uval{5x+\hex{4f19}}$ \\
1 & case 0 & $\uval{3x+\hex{3a75}}$ \\
2 & case 1 & $\uval{(x\oplus\hex{643})+(x\gg1)}$ \\
3 & case 3 & $\uval{4x}\oplus\hex{7b03}$ \\
4 & case 6 & $\uval{(x\oplus\hex{d854})+(x\gg1)}$ \\
5 & case 7 & $\uval{4x}\oplus\hex{7d53}$ \\
6 & case 4 & $\uval{5x+\hex{26cb4}}$ \\
7 & case 2 & $\uval{5x+\hex{4385a}}$ \\
\bottomrule\end{tabular}\end{minipage}


代码先输出 i，再用 switch 计算第二项。表格中的常量来自个人炸弹的实际指令，保存的是题目规则，不是几个 seed 的现成答案。例如右移与异或相加的分支，需要保留原来的 x；不能先覆盖 x，再拿改变后的值参与另一半运算。

执行第 3 关公开测试，10 个 seed 全部通过。

\reportimage{images/phase3-public.png}{0.34}{图 6：第 3 关公开测试通过 10/10。}

\reportimage{images/phase3-seed8.png}{0.45}{图 7：第 3 关使用非公开 seed 8 运行，输出 \texttt{Phase 3 defused.}}


\Needspace{222pt}\subsection*{第 4 关：递归查找与逐层返回}


本关输出一个整数。入口设置目标、密钥和初始区间：
\[
  x=c\band127,\qquad
  k=\uval{c\gg7}\oplus\uval{c\gg32},\qquad [l,r]=[0,127].
\]
helper 每次取区间中点，再选择包含目标的那一半。

\noindent\begin{minipage}{\linewidth}
\begin{lstlisting}[language={}]
402328: sub %esi,%edi
40232a: shr $1,%edi
40232c: add %esi,%edi
40232e: cmp %eax,%edi
402330: je 402380 <p4_walk+0x60>
402337: lea (%rdi,%rcx,1),%ebx
40233e: cmp %edi,%eax
402340: jb 402360 <p4_walk+0x40>
402347: call 402320 <p4_walk>
40234c: lea (%rax,%rax,4),%eax
402356: xor $0x17,%eax
402365: call 402320 <p4_walk>
402374: xor $0x1f,%eax
402380: lea 0x68(%rdi),%eax
\end{lstlisting}
\end{minipage}


设当前区间的中点为
\[
  m=l+\left\lfloor\frac{r-l}{2}\right\rfloor.
\]
将递归函数记为 $F(x,l,r;k)$，三种返回情况分别为
\[
F(x,l,r;k)=
\begin{cases}
  \uval{m+\hex{68}}\oplus k, & x=m,\\[3pt]
  \uval{5F(x,l,m-1;k)}\oplus\uval{m+k}\oplus\hex{1f}, & x<m,\\[3pt]
  \uval{5F(x,m+1,r;k)}\oplus\uval{m+k}\oplus\hex{17}, & x>m.
\end{cases}
\]
本关口令为 $F(x,0,127;k)$。其中左、右分支使用的常量分别是 $\hex{1f}$ 和 $\hex{17}$，递归过程中 $k$ 始终不变。

因此只找到目标位置还不够，必须按调用栈由内向外合并返回值。实现用 \texttt{walk4} 保留相同递归结构，左右分支分别写清参数与常量。目标原本就在合法区间内，且每次区间缩小，因此会到达命中条件；不需要增加与实际程序无关的空区间返回规则。

执行第 4 关公开测试，结果为 10/10。

\reportimage{images/phase4-public.png}{0.34}{图 8：第 4 关公开测试通过 10/10。}

\reportimage{images/phase4-seed8.png}{0.45}{图 9：第 4 关使用非公开 seed 8 运行，输出 \texttt{Phase 4 defused.}}


\Needspace{222pt}\subsection*{第 5 关：字段偏移与链表遍历}


本关要求两个整数，分别是累计结果和最后访问节点的 tag。节点区域起点是 \texttt{0x405ba0}，相邻节点相距 16 字节：偏移 0 为 value，偏移 4 为 tag，偏移 8 为 next 指针。片段中的 \texttt{0x405b60} 是八项头节点索引表，不是节点区本身。

\noindent\begin{minipage}{\linewidth}
\begin{lstlisting}[language={}]
40241d: and $0x7,%eax
402420: shr $0x3,%rcx
402424: mov 0x405b60(,%rax,8),%rdx
402433: shr $0x6,%rax
402437: xor %edi,%eax
402439: mov 0x4(%rdx),%esi
40243c: mov %ecx,%edi
40243e: xor $0xbf49,%eax
402443: lea (%rax,%rax,2),%eax
402446: add (%rdx),%eax
40244b: xor %esi,%eax
40244d: and $0x3,%edi
402457: mov 0x8(%rdx),%rdx
402461: mov 0x4(%rdx),%esi
402466: xor %esi,%eax
\end{lstlisting}
\end{minipage}


将头节点索引表记为
\[
  H=(11,9,3,5,6,10,1,0).
\]
令 $p_j$ 表示第 $j$ 次访问的节点索引，$V_p$、$T_p$、$N_p$ 分别表示节点 $p$ 的 value、tag 和 next 字段。初始节点、访问次数和累计值为
\[
\begin{aligned}
p_0&=H_{c\band7},\\
n&=1+((c\gg3)\band3),\\
s_0&=\uval{c\gg6}\oplus\uval{c\gg32}\oplus\hex{bf49}.
\end{aligned}
\]
每访问一个节点，就更新累计值并沿 next 移动：
\[
\begin{aligned}
s_{j+1}&=\uval{3s_j+V_{p_j}}\oplus T_{p_j},\\
p_{j+1}&=N_{p_j},\qquad j=0,\ldots,n-1.
\end{aligned}
\]
两个输出依次为 $s_n$ 和 $T_{p_{n-1}}$，即累计结果与最后一个已访问节点的 tag。

头节点已经在循环前处理了一次，之后最多再走三次，所以 $n\in\{1,2,3,4\}$。代码用结构体数组保存个人炸弹的数据，用索引表达 next，避免将程序中的固定地址直接当成自己进程里的指针。

例如物理索引 11 的 next 是 9，接着是 3、5；它不是按节点在文件里的物理位置顺序移动。当前头表和步数范围保证访问完成前不会遇到 NULL。完整的十二个节点 value、tag、next 保存在 \texttt{defuser.c} 的 nodes 数组中。

\reportimage{images/phase5-public.png}{0.34}{图 10：第 5 关公开测试通过 10/10。}

\reportimage{images/phase5-seed8.png}{0.45}{图 11：第 5 关使用非公开 seed 8 运行，输出 \texttt{Phase 5 defused.}}


\Needspace{200pt}\subsection*{第 6 关：函数指针表与两次间接调用}


本关输出一个整数。入口把 challenge 交给 \texttt{p6\_dispatch}，后者从函数指针表中选择两次操作，第一次结果还要经过异或才能作为第二次的输入。

\noindent\begin{minipage}{\linewidth}
\begin{lstlisting}[language={}]
40256c: shr $0x6,%rdi
402573: shr $0x20,%rax
402577: xor %eax,%edi
402579: mov %rbx,%rax
402580: and $0x38,%ebx
402583: and $0x7,%eax
402586: call *0x405e00(,%rax,8)
40258d: xor $0xcb24,%eax
402592: mov %eax,%edi
402594: call *0x405e00(%rbx)
40259f: add $0x953e,%eax
\end{lstlisting}
\end{minipage}


函数表位于 \texttt{0x405e00}，其映射关系及两次索引为
\[
\begin{aligned}
P&=(1,3,2,2,4,3,0,5),\\
i_0&=c\band7,\qquad i_1=(c\gg3)\band7.
\end{aligned}
\]
汇编保存的 $c\band\hex{38}$ 是第二个表项的字节偏移，与 $8i_1$ 等价。将六种操作记为 $f_0,\ldots,f_5$，对应关系如下。

\par\addvspace{5pt}\noindent\begin{minipage}{\linewidth}\centering\small\renewcommand{\arraystretch}{1.28}
\begin{tabular}{>{\raggedright\arraybackslash}p{0.140\dimexpr\linewidth-6\tabcolsep\relax}>{\raggedright\arraybackslash}p{0.430\dimexpr\linewidth-6\tabcolsep\relax}>{\raggedright\arraybackslash}p{0.430\dimexpr\linewidth-6\tabcolsep\relax}}
\toprule
\textbf{操作编号} & \textbf{计算关系} & \textbf{观察重点} \\
\midrule
0 & $\operatorname{ror}_{32}(x,7)$ & 右侧移出的位回到左侧 \\
1 & $\uval{2x+\hex{b13c}}\oplus\uval{7x}$ & 两个表达式都使用原来的 $x$ \\
2 & $\uval{(x\oplus(x\gg5))+\hex{f09a}}$ & 异或后再加常量 \\
3 & $\uval{x+\hex{484b}}\oplus(x\gg5)$ & 加法与右移分别计算 \\
4 & $\operatorname{rol}_{32}(x,1)\oplus\hex{7481}$ & 左侧移出的位回到右侧 \\
5 & $\uval{(x\oplus\hex{85d4})+4x}$ & 先异或，再与原值四倍相加 \\
\bottomrule\end{tabular}\end{minipage}


两次间接调用与最后的加法可写成
\[
\begin{aligned}
x_0&=\uval{c\gg6}\oplus\uval{c\gg32},\\
x_1&=f_{P_{i_0}}(x_0)\oplus\hex{cb24},\\
v&=\uval{f_{P_{i_1}}(x_1)+\hex{953e}}.
\end{aligned}
\]
实现用 \texttt{phase6\_op} 的 switch 表达六种操作，再由 ops 数组恢复实际表序。这里 $\operatorname{ror}_{32}$、$\operatorname{rol}_{32}$ 分别表示 32 位循环右移、循环左移；循环移位用两次方向相反的移位和按位或实现，普通移位不能代替循环移位。

执行第 6 关公开测试，10 个 seed 全部通过。

\reportimage{images/phase6-public.png}{0.34}{图 12：第 6 关公开测试通过 10/10。}

\reportimage{images/phase6-seed8.png}{0.45}{图 13：第 6 关使用非公开 seed 8 运行，输出 \texttt{Phase 6 defused.}}


\Needspace{222pt}\subsection*{Honor：本地对象布局与第七行文本}


Honor 接收 challenge 和由前六行计算出的 context。入口不只是读取三个整数，还会先检查整行长度、复制文本、检查对象内的相邻字段，再从后续位置解析整数。

\noindent\begin{minipage}{\linewidth}
\begin{lstlisting}[language={}]
40272b: lea -0xd0(%rbp),%r12
402751: sub $0x24,%rax
402755: cmp $0x7f,%rax
402759: jbe 402770 <secret_gate+0x60>
402773: mov %r12,%rdi
402776: call 4026f0 <secret_copy>
40277b: cmpl $0x4e45504f,-0xb0(%rbp)
40278d: lea -0xac(%rbp),%rdi
402794: call 402630 <secret_phase>
40268b: cmp $0x3,%eax
402614: lea 0xfca2(%rdi,%rdi,4),%eax
40261b: xor $0x525b,%eax
\end{lstlisting}
\end{minipage}


复制起点为 \texttt{rbp-0xd0}，被比较字段在 \texttt{rbp-0xb0}，相差 32 字节。按小端字节序，\texttt{0x4e45504f} 对应 \texttt{OPEN}。随后送入 \texttt{secret\_phase} 的地址又向后移动 4 字节，所以文本结构为“32 个填充字符 + OPEN + 三个整数”。长度检查允许 36 到 163 字节；当前构造的文本在此范围内。

设 $h$ 为 context，并令 $k$ 对应代码中的 kind。三个整数按以下关系生成：
\[
\begin{aligned}
k&=h\band3,\\
a&=\uval{c}\oplus\uval{c\gg32},\\
b&=\uval{h}\oplus\uval{h\gg32},\\
x&=\uval{5(a\oplus b)+\hex{fca2}}\oplus\hex{525b},\\
y&=x\oplus\uval{k\cdot\hex{fca2}}\oplus\uval{h\gg16},\\
z&=y\oplus\uval{b+k+\hex{525b}+(a\gg7)}.
\end{aligned}
\]


整数顺序为 $k$、$y$、$z$。context 必须使用当前六行输出对应的值，不能固定为某次运行的常数；现有驱动已完成这一步并把它作为实参传入。实现用 \texttt{snprintf} 构造第七行，并检查返回长度，保证输出缓冲区没有截断。

执行 Honor 公开测试，10 个 seed 全部通过。该测试先验证前六行，再验证对应的第七行，因此也检查了前六行与 context 的配合。

\reportimage{images/honor-public.png}{0.34}{图 14：Honor 公开测试通过 10/10。}

再使用公开列表之外的 seed \texttt{8} 验证 Honor。先输出前六关答案，再输出对应的第七行文本，将两部分一并交给相同 seed 的炸弹：

\begin{lstlisting}[language={}]
{
    ./defuser --seed 8
    ./defuser --seed 8 --secret
} | ./bomb --seed 8
\end{lstlisting}

图 15 记录了这次运行的命令与启动画面；成功结果另经真实 Ubuntu 终端复核，日志末尾为：

\begin{lstlisting}[language={}]
Congratulations! You've defused the bomb.
Secret phase defused.
\end{lstlisting}

这说明非公开 seed 下，前六关与 Honor 的第七行能够连续通过。完整复核输出随报告保存在 \href{run:logs/honor-seed8.log}{\texttt{logs/honor-seed8.log}}。

\reportimage{images/honor-seed8.png}{0.43}{图 15：非公开 seed 8 的完整运行命令与启动画面。}



\clearpage
\Needspace{6\baselineskip}\section*{三、验证过程与结果}


代码在真实 Ubuntu 中按课程原有编译选项完成构建。运行单关公开测试和 \texttt{public\_test.py} 总测试，常规六关通过 60/60，Honor 通过 10/10，结果见图 1 至图 16。

本次截图记录了非公开 seed \texttt{8} 在第 1 至第 6 关的单关运行结果（图 3、图 5、图 7、图 9、图 11、图 13）。公开 seed 列表见 \texttt{tests/public.json}，\texttt{8} 不在其中，这些试跑属于公开列表之外的额外验证。Honor 同样使用非公开 seed 8 完成验证，命令见图 15，通过结果见上一节的终端日志摘录。

\par\addvspace{5pt}\noindent\begin{minipage}{\linewidth}\centering\small\renewcommand{\arraystretch}{1.28}
\begin{tabular}{>{\raggedright\arraybackslash}p{0.130\dimexpr\linewidth-6\tabcolsep\relax}>{\raggedright\arraybackslash}p{0.670\dimexpr\linewidth-6\tabcolsep\relax}>{\raggedright\arraybackslash}p{0.200\dimexpr\linewidth-6\tabcolsep\relax}}
\toprule
\textbf{关卡} & \textbf{验证记录} & \textbf{实验截图} \\
\midrule
第 1 关 & 公开 10/10；seed 42、seed 8 复跑 & 图 1、图 2、图 3 \\
第 2 关 & 公开 10/10；非公开 seed 8 复跑 & 图 4、图 5 \\
第 3 关 & 公开 10/10；非公开 seed 8 复跑 & 图 6、图 7 \\
第 4 关 & 公开 10/10；非公开 seed 8 复跑 & 图 8、图 9 \\
第 5 关 & 公开 10/10；非公开 seed 8 复跑 & 图 10、图 11 \\
第 6 关 & 公开 10/10；非公开 seed 8 复跑 & 图 12、图 13 \\
Honor & 公开 10/10；非公开 seed 8 通过 & 图 14、图 15 \\
总测试 & 常规六关 60/60；Honor 10/10 & 图 16 \\
\bottomrule\end{tabular}\end{minipage}

\Needspace{6\baselineskip}\section*{四、关键问题与理解}


分析的第一步是把输入数与检查条件对应起来，并跟踪 helper 调用前后的参数和返回值。第 3、6 关的表项可能经过重排，因此表索引不等于函数编号；第 5 关中，节点的物理排列也不等于 next 所决定的访问顺序。第 4 关的返回值是在递归返回时逐层合成，而不是只由最终目标决定。

位宽同样需要单独检查。64 位 challenge 的移位必须发生在截成 32 位之前，循环移位与普通移位不同，32 位加乘也不能随意改用有符号整数。修改代码后应重新编译，再使用相同命令运行，以免误把旧可执行文件的结果当作新代码的结果。

\Needspace{6\baselineskip}\section*{五、资料与辅助工具使用}

报告使用 \LaTeX{}（XeLaTeX，TeX Live 2025）排版，图片来自本人运行时的终端截图。写作过程中使用 AI 助手协助 LaTeX 排版、汇编核对与最终终端复核。参考资料为课程仓库的 \texttt{README.md}、\texttt{tutor/} 与 \texttt{teaching/} 目录。

\clearpage
\Needspace{6\baselineskip}\section*{附：公开测试总汇总}


\reportimage{images/public-summary.png}{0.70}{图 16：总测试输出，常规六关通过 60/60，Honor 通过 10/10。}

\par\bigskip
\Needspace{6\baselineskip}
希望我能在看到 \texttt{boom!} 和 \texttt{PASS} 之后，想起某年某月某天为这个实验伤心和开心了很久之后的自己。

感谢远道而来的实验和助教的用心。疲惫之后，记得跟彼此的自己说一句，早安。

\end{document}
